\section{HECTOR Framework}
%\subsection{MLIR Infrastructure}
% \hl{let us briefly discuss this in flow.}

% Hector is built on MLIR\cite{MLIR}, a novel compiler infrastructure providing strong scalability and modularity, and is particularly suitable for multi-level IR transformation. The \emph{Dialect} in MLIR is a set of operations, attributes and types, which can be regarded as a flexible IR. All the IRs in MLIR obey Static Single Assignment (SSA) form \cite{StaticSingleAssignment} and explicit type system. MLIR provides several built-in Dialects, such as \texttt{SCF} and \texttt{Standard}. \texttt{SCF} describes static control flow in a higher-level abstraction keeping branch and loop structure, which is helpful for analyzing and optimizing. \texttt{Standard} is a collection of basic operation such as addition, comparison and function call. These dialects can be used to describe a program at the same time, which is also the feature provided by MLIR, called progressive lowering. The program described in \texttt{SCF} and \texttt{Standard} has similar functionality and representation with high-level language.

%\subsection{Overall flow}

% \begin{figure}[t]
%   \centering
%   \includegraphics[width=\linewidth]{figure/overview.pdf}
%   \caption{Overview of HLS flow with Hector framework}
%   \label{fig:flow} \vspace{-0.3cm}
% \end{figure}

% \autoref{fig:flow} presents the overall flow of Hector. The framework is built upon MLIR  \cite{MLIR} ecosystem, a highly customizable compiler infrastructure, which greatly facilitates the implementation of various IRs and transformations among them. The input to Hector is standard C or C++ program. We adopt open-source tools Clang \cite{lattner2008llvm} and Polygeist \cite{polygeist}, which lower C/C++ description to MLIR system in a two-step manner. Hector then performs an efficient algorithm (as detailed in \autoref{sec:algorithm}) to partition the design into submodules, synthesizing them progressively in a hybrid manner, and generate the Chisel description at the register transfer level (RTL). 

%Our framework not only supports versatile scheduling techniques, but also provides lots of customization and extension opportunities thanks to the implementation of algorithms upon MLIR infrastructure. 

%Hector applies compilation optimizations, such as constant propagation, on \texttt{scf} dialect, an existing software IR with high-level control flow semantics in MLIR system. The optimized \texttt{scf} code is transformed into a raw \tor description through operation rewriting and topology initiation. 

% Hector is essentially a two-level IR system equipped with a set of optimization and transformation passes. We design \tor IR at a high level and \hec IR at a low level, as detailed in \autoref{sec:ir}. \tor (Topology Rail) IR features a directed graph, where high-level operations are bound on nodes or edges. Compared with pure software IRs, such as LLVM IR \cite{lattner2008llvm}, \tor is much more capable for hardware design because the graph carries rich hardware semantics. \hec (Hybrid Elastic Component) IR is designed at the near-RTL level. It describes each component in either finite-state machine (FSM) or pipelining manner, and integrates them using elastic units. The hybrid nature of \hec greatly facilitates the description of hardware generated by versatile HLS methodologies, while existing hardware IRs either support only a specific type of hardware (\cite{Calyx_Rachit} only uses FSM to describe hardware control, and \cite{uIR_Sharifian} only describes dataflow circuits, for example) or \cite{LLHD_Schuiki} describes digital circuits at a much lower level not suitable for HLS flow.

%and bridges the semantic gap between Von Neumann programming model and fine-grained parallel hardware. 


% After execution, Hector would deliver the produced RTL description with appropriate IP configurations to the commercial software Vivado for logic synthesis and physical design, which yields bitstream results together with analysis reports.

% \autoref{fig:flow} depicts how a high-level program is transformed to hardware in Hector. Hector first compiles C/C++ programs to clang AST, which Polygeist\cite{polygeist} traverses to create \texttt{SCF} + \texttt{Standard}. Some essential optimizations like constant propagation are applied, and a raw \dsir is built by rewriting operations and initializing the schedule topology. Hector can construct a hybrid scheduling circuit and then adopt a static scheduling techniques on static modules, both of which can be simply implemented on the topology structure of \dsir. The schedule results are then lowered to \fsm composed of different components using typical static and dynamic HLS synthesis processes. \fsm components are directly transformed to Chisel\cite{chisel} programs that describes the hardware using some pre-defined Chisel templates. The simulation provided by Chisel can be used to validate the functionality and performance, and Chisel can generate results written in Verilog which fits in logic synthesis and physical design.

% We design two a two-level IR, \dsir and \fsm, which are hardware descriptions at different abstraction level. \dsir{} provides a hardware description IR with a well-designed type system and operations, which represents static and dynamic scheduled compound design in a hierarchical description. \fsm{} reduces the gap between two scheduling strategies, and is also a hardware-like structure which can be easily translated to hardware implementation like HDL. These two dialects are introduced in \autoref{sec:dialect}. 

% Our proposed flow is built on these two \emph{Dialects} and MLIR build-in high-level dialects. The input dialects includes \emph{SCF} and \emph{Std}. These dialects together describe a high-level program, which reserves the structures of branch and loop. Compared with low-level representation CDFG, the structures of branch and loop are friendly to understand and scheduling. Furthermore, \emph{Affine} defines a collection of affine operation, which can make use of polyhedral model analysis to achieve loop transformation and memory partition. Optimization directives like loop unroll and memory partition can also be inserted into this dialect. Then, static scheduling is applied to the high-level program, where SDC algorithm is provided in this flow. This framework can also support other scheduling algorithm that is implemented by users. Static scheduling process will generate fully static modules represented in \dsir{}. Then, static modules are transformed to mixed modules by the a partition algorithm, where mixed modules are the combination of static and dynamic modules. The mixed modules are then split into static and dynamic modules, both of which have further transformation passes. After that, both of modules are lowered to \fsm{}. Having done some low level transformations in \fsm{}, the programs are then converted to Chisel\cite{chisel}. Detailed design and the reason why partition algorithm is applied to the static modules are introduced in the following sections.
